\section{Combining the local constructions}
\label{sec-global}

Fix an integer $r\ge6$ divisible by six, and put $k=r-1$ and $d=r-3$.
We use the labelled frame family of the preceding section, with
\[
 \Delta=64r^8,\qquad
 \eta=\frac{r!}{2d!(\Delta+1)},\qquad
 \beta=(r-2)\eta.
\]
At each sufficiently large background size $M$, select exactly
$\lfloor\eta M^d\rfloor$ labelled frames whose allowed sets of cyclic
rows are pairwise disjoint. The selection is fixed before the modifications
described below. Define
\begin{equation}\label{eq-global-parameters}
 A=\left\lceil\frac{r^3}{\beta}\right\rceil,\quad
 Q_{
 \mathrm{grid}}=rA,\quad W=r^2A,\quad H=rA^2,\quad
 G=W-r,\quad \rho=\frac{A\beta}{r},\quad
 E=10^{12}r^{22}W^4.
\end{equation}
The letter $Q_{\mathrm{grid}}$ denotes a grid width. It is unrelated to
the row charge used in the assembly section.

\subsection{Sets of active letters and fixed background letters}

Take the $W$ control letters to be the pairs $(i,j)$ with
$0\le i<r$ and $j\in\mathbb Z/Q_{\mathrm{grid}}\mathbb Z$.
For $0\le a<A$ and $b\in\mathbb Z/Q_{\mathrm{grid}}\mathbb Z$, set
\[
 T_{a,b}=\{(i,ai+b):0\le i<r\}.
\]
The members of $T_{a,b}$ are the \emph{active letters} of that
construction. The other $G$ control letters are its \emph{guards}; these
letters stay in the background circle and are never chosen as mobile
letters. For example, when $r=3$ and $A=2$, the blocks $T_{0,0}$ and
$T_{1,0}$ in a grid with six columns share just $(0,0)$.

There are $H$ blocks, each with $r$ letters. Distinct blocks intersect
in at most one letter. Indeed, two common letters in distinct rows $i,j$
would imply
\[
 Q_{\mathrm{grid}}\mid(a-a')(i-j),\qquad
 |(a-a')(i-j)|\le(A-1)(r-1)<Q_{\mathrm{grid}}.
\]
The product must therefore be zero. Since $i\ne j$, this gives $a=a'$,
and either common letter then gives $b=b'$. This argument uses integers
and remains valid when $Q_{\mathrm{grid}}$ is composite. Each letter
belongs to $A$ blocks. It follows that a block meets exactly
$I=r(A-1)\le W/r$ other blocks: each of its letters lies in $A-1$ other
blocks, and no block is counted at two letters.

Fix a cyclic order of $m$ old letters. Insert each guard in a different
old cut, immediately before the old letter at that cut. A \emph{guard
background} is the resulting cyclic order of $M=m+G$ letters. There are
exactly
\[
 (m)_G=m(m-1)\cdots(m-G+1)
\]
guard backgrounds for each block. Distinct cuts give no ordering
multiplicity. For $m\ge G$,
\begin{equation}\label{eq-guard-count}
 1-\frac{G(G-1)}{2m}\le\frac{(m)_G}{m^G}\le1.
\end{equation}
The lower bound follows by applying
$\prod_i(1-u_i)\ge1-\sum_i u_i$ to $u_i=i/m$.

A \emph{candidate} consists of an active block, a guard background, and
one selected labelled frame. Candidates on the same block but different
backgrounds cannot share a row: deleting their active letters would recover
the same background from that row. Candidates on the same background are
already separated by the selected frame family. Thus only candidates on
different blocks can collide.

\subsection{A finite list of necessary collision equations}

Anchor every cyclic background at the same old letter. A local background
cut $y\in\{0,\ldots,M-1\}$ projects to an original cut
$x\in\{0,\ldots,m-1\}$. Counting the guards passed before $y$ gives
\[
 y\equiv x+e\pmod m,\qquad 0\le e\le G.
\]
This also covers the guard immediately before the anchor. In that case
the local cut has the form $m+e$ and the original cut is zero.

Every stationary frame satisfies a relation whose coefficient multiset is
\[
 (d-1,1,\underbrace{-1,\ldots,-1}_{d\text{ entries}}).
\]
The coefficients sum to zero, and their absolute values sum to $2d$.
Including the binary cut phases and reducing the local relation modulo
$m$ shows that any foreign row must satisfy
\begin{equation}\label{eq-foreign-relation}
 \sum_u c_u x_u\equiv b\pmod m,\qquad
 |b|\le C_1=(4d+1)G+2d\le4rW.
\end{equation}
For completeness, choose the canonical local cuts in $[0,M)$.
Their weighted sum differs from the binary phase sum by an integer
multiple of $M$, whose multiplier has absolute value at most $2d+1$.
Replacing $M$ by $m+G$ and each local cut by $x_u+e_u$ gives the bound
$(2d+1)G+2dG+2d$ displayed above.
There are $k(k-1)$ assignments of this coefficient multiset to the
stationary physical labels, since one chooses the label with coefficient
$d-1$ and then the label with coefficient $1$.

First delete a candidate if a relation \eqref{eq-foreign-relation} holds
for a possible foreign stationary set containing none of its active
letters. Such a relation has a nonzero coefficient on a guard variable.
There are at most
\[
 J_1\le Hr\,k(k-1)(2C_1+1)\le9rW^3
\]
relations to consider. A nonzero coefficient has absolute value at most
$r$. After all other guard variables are fixed, its congruence has at most
$r$ solutions modulo $m$, because the number of roots of $ax=b$ is at most
$\gcd(a,m)\le|a|$. Thus each equation excludes at most $rm^{G-1}$
guard tuples. If $m\ge G^2$, \eqref{eq-guard-count} gives
$(m)_G\ge m^G/2$. This first deletion removes at most the fraction
\begin{equation}\label{eq-guard-loss}
 \frac{18r^2W^3}{m}.
\end{equation}
The selected frame count cancels from this calculation.

Suppose two surviving candidates on different blocks share a cyclic row.
Their blocks must have a common letter $v$, and $v$ must be stationary in
both rows. If the blocks were disjoint, all foreign stationary letters
would be guards. If $v$ were mobile in one row, its stationary letters
would all be guards for the other candidate. Either case contradicts the
first deletion, applied to the appropriate candidate.

If $v$ occupies stationary slot $j$ during batch $t$, then its local cut
is $(q_j+t+\epsilon_j)\bmod M$. Here $0\le t<M$ is the canonical
batch number. The local wrap is zero or one, so its original cut obeys
\[
 x_v\equiv q_j+t+e'\pmod m,\qquad -2G\le e'\le1.
\]
Substituting in \eqref{eq-foreign-relation} gives the following finite
list of possible events:
\begin{equation}\label{eq-event}
 a(t+q_j)+\sum_u b_ux_u\equiv b_1\pmod m,
 \qquad |b_1|\le C_2=(6d-1)G+3d-1\le6rW.
\end{equation}
An \emph{event} forbids the pair $(v,j)$ for that entire batch.
The list ranges over every meeting foreign block, every foreign mobile
letter, every coefficient assignment, every candidate slot $j$, and every
integer $b_1$ in the stated interval. It therefore includes both the
original rows and every permitted later modification.

There are at most
\[
 J_{\mathrm{event}}\le I k^3(k-1)(2C_2+1)\le13r^4W^2
\]
equation types. For a fixed type there are at most $2|a|$ actual roots
$0\le t<M$, provided $M<2m$. For a fixed shared label, summing $|a|$
over its coefficient assignments gives $2d(d+1)$. Hence the total number
of events, with repetitions allowed, is at most
\begin{equation}\label{eq-event-count}
 K\le4Ik^2d(d+1)(2C_2+1)\le52r^4W^2.
\end{equation}

\subsection{Separating events without a distribution assumption}

Put $L=2r+1$ and $S=2L$. Delete a candidate if it has two events for
different forbidden pairs at circular distance at most $S$ on the
$M$-batch circle. Eliminating the two event times gives
\begin{equation}\label{eq-proximity}
 aa'(q_j-q_l)+\sum_u(a'b_u-ab'_u)x_u
 \equiv a'b_1-ab'_1+aa'(\delta+\sigma G)\pmod m,
\end{equation}
where $|\delta|\le S$ and $\sigma\in\{-1,0,1\}$.
The term $\sigma G$ accounts for a wrap on the $M$-circle. It must be
retained even though the congruence is modulo $m$.
There are at most
\[
 J_2\le3(2S+1)J_{\mathrm{event}}^2\le4563r^9W^4
\]
such equations.

Every equation \eqref{eq-proximity} has a nonzero variable coefficient.
Here are all possibilities. If the foreign blocks differ, each has
$r-2$ guard letters in its stationary set, and the two blocks share at
most one letter. A private guard then has a nonzero coefficient. If the
foreign block is the same but its mobile letters differ, a guard present
in only one stationary set has a nonzero coefficient. If the block and
mobile agree, either a guard coefficient $a'b_u-ab'_u$ is nonzero or the
two full coefficient vectors are proportional. In the latter case their
equal multisets force the ratio to be $1$. Indeed their squared norms
force the ratio to be $1$ or $-1$, and $-(d-1)$ does not occur in the
multiset when $d\ge3$. Thus the vectors are equal. The forbidden pairs
then differ only in their slots, so the remaining form is
$a^2(q_j-q_l)$ with $j\ne l$. On a fixed position of the sum point,
the cuts have the distinct coefficient vectors
$0,e_1,\ldots,e_d,\boldsymbol{1}$. Their difference is nonzero.
All coefficients have absolute value at most $3r^2$.

We count the guard-variable and guard-free equations separately. If a
guard coefficient is nonzero, fix the selected labelled frame and all
other guards. There are at most $3r^2$ roots modulo $m$. Dividing by
$(m)_G\ge m^G/2$ bounds the deleted fraction by $6r^2/m$ per equation,
uniformly in the selected frame.

If all guard coefficients vanish, partition the sorted frame parameters
according to the position of their sum point and the integer
\[
 \ell=\left\lfloor\frac{\sum_i z_i}{M}\right\rfloor.
\]
There are at most $d(k-1)\le r^2$ choices. Within such a part,
$q_j=v_j\cdot z-\mathbf{1}_{j=a}\ell M$. The winding term changes only
the constant in \eqref{eq-proximity}. Fix a coordinate with nonzero
coefficient. It ranges in an interval shorter than $2m$, so it has at
most $6r^2$ roots. Deleting that coordinate leaves a strictly increasing
list of $d-1$ integers, with at most $M^{d-1}/(d-1)!$ possibilities.
Including all $r!$ initial assignments and dividing by
$\lfloor\eta M^d\rfloor\ge\eta M^d/2$ bounds the deleted fraction,
per part and per equation, by
\[
 \frac{12r^2r!}{\eta(d-1)!M}
 =\frac{24r^2d(\Delta+1)}{M}.
\]
This counts in the entire parameter set before restricting to the selected
frames. It requires no uniform distribution property of that selection.
Since $\Delta+1\le65r^8$ and $M\ge m$, both cases together are bounded by
\begin{equation}\label{eq-proximity-loss}
 \frac{7\,118\,280r^{22}W^4}{m}.
\end{equation}

\subsection{Modifying every surviving candidate}

Around every event take the closed interval of radius $L$ on the batch
circle, and merge intersecting intervals. Each connected component has
one forbidden pair, since intervals for different pairs cannot intersect.
If there are no events, keep the normal construction. Otherwise,
$M>(2L+1)K$ ensures a gap longer than $2L+1$ between consecutive event
times. Cut in this gap and regard the remaining components as ordinary
intervals. The cut changes only where these intervals are written; the
canonical times in \eqref{eq-event} remain unchanged.

For a component whose first and last events are $f,e$, let $(v,j)$ be
its forbidden pair. Choose the normal boundary $s\in[f-r+1,f]$ at
which $v$ is mobile; the normal role action is a cycle on the $r$ roles.
Set
\[
 P=2r\left\lceil\frac{e+1-s}{2r}\right\rceil.
\]
Replace the next $P$ normal batches by $P/2$ pairs of a permitted order
and its reverse, using the order whose first target is different from
$j$. Follow them by one normal cleanup batch. The local endpoint theorem
shows that these $P+1$ batches end at exactly the normal literal state.
Throughout the first $P$ batches, the letter $v$ is either mobile or in
the first target slot. Thus it never occupies $j$ during any event batch.

The inequalities
\[
 f-r+1\le s\le f,\qquad
 e+1\le s+P\le e+2r,\qquad s+P+1\le e+2r+1
\]
put the entire replacement inside its merged interval. Different
replacements are disjoint. Each restores its normal endpoint, so they can
all be performed on the same cyclic walk. The joint-state containment
proved earlier keeps every modified row inside the originally selected
allowed set. If two modified candidates met, their unique shared letter
would be stationary in both rows and would produce an event
\eqref{eq-event}. The modification prevents that occupancy. This excludes
collisions between two modifications as well as collisions with normal
rows.

Regard a batch starting at integer time $t$ as the unit interval $[t,t+1]$.
The inequalities above put every changed batch inside the interval union.
That union has total length at most $2LK\le5rK$, so at most $5rK$
batches change. A complete batch has
at most $k(M+k)$ source rows. Comparing each changed batch with its normal
version therefore changes the total row count by at most $10r^2KM$ when
$M\ge r$. The short charge stays exactly $2(r-1)M$. There is also an
untouched normal batch. Its full rows follow the literal successor $F$,
which is the additional hypothesis needed in the connector construction.

\subsection{Uniform counts}

Let $B$ be the total number of source rows in the surviving modified
contours, let $q$ be their total short charge, and let $c_{\mathrm{adv}}$
be the number of circles in the actual connector lists constructed in the
assembly section. Put
$D_{m,W}=\prod_{i=0}^{W-1}(m+i)$ and $c_r=2(r-1)/(r-2)$.

\begin{proposition}\label{prop-global-counts}
For every multiple $m$ of $r$ with $m\ge4E$, the construction has literal
closure, pairwise distinct source cyclic classes above the chosen old
cyclic order, and the safe connector lists of the assembly section. Its counts satisfy
\[
 \frac{B}{D_{m,W}}=\frac{\rho W}{m}(1+\theta),\qquad
 \frac{q+c_{\mathrm{adv}}}{D_{m,W}}
   =\frac{c_r\rho W}{m^2}(1+\epsilon),\qquad
 |\theta|,|\epsilon|\le\frac Em,
\]
and $c_{\mathrm{adv}}\le(r-1)q/m$.
\end{proposition}

\begin{proof}
We first verify the size conditions used in the construction. The ceiling
defining $A$ gives
\[
 \eta^{-1}=\frac{r-2}{\beta}\le\frac{r}{\beta}
 \le\frac{W}{r^4}\le W.
\]
Also $W\ge r^2\ge1$ and $G\le W$. The definition of $E$ implies
$E\ge W^2$, $E\ge rW$, and $4E\ge38r^4$. Hence
\[
 M\ge m\ge38r^4,\qquad
 m\ge4W^2\ge G^2,\qquad
 G<m,\qquad M=m+G<2m.
\]
Since $d\ge3$ and $M\ge m\ge4W$, we have
$\eta M^d\ge M/W\ge4$. Moreover,
$rG/m\le rW/(4E)\le1/4$. Finally, the event bound gives
\[
 (4r+3)K\le260r^5W^2<4\cdot10^{12}r^{22}W^4\le m\le M.
\]
The strict middle inequality follows from $r,W\ge1$ and
$260<4\cdot10^{12}$. These prove all the required size conditions.
Also $r\mid M$, since $r$ divides $m$ and $W$. The local closure and
distinctness results and the preceding modifications consequently apply.
Every source class is above the chosen old cyclic order. Deleting all
controls first removes the guards from the guard background. A full move
only rotates the remaining old cyclic order. A positive short update
moves a control past one background letter, so it does not change the
relative cyclic order of the old letters. This remains true during every
replacement. Deleting all controls from any resulting row therefore
recovers the chosen old class, as required by
Lemma~\ref{lem-assembly-universe}.

Each modified contour has
\[
 N=(r-2)M^2+[r(r-1)+1]M+\delta_N,
 \quad |\delta_N|\le10r^2KM,\quad q_0=2(r-1)M.
\]
Let $h=M+r-1$ be the source port modulus and let $g$ be its greatest
common divisor with the monodromy. The local result gives $1\le g\le r-1$.
Write $k=r-1$. The untouched normal batch has $k$ short rows and
$(r-2)M+r(r-1)+1$ rows in total. One of its $k$ consecutive full runs
therefore has length
\[
 R\ge\frac{(k-1)M+k^2+1}{k}.
\]
Its bases follow the literal successor $F$, and the local construction
gives $R\le h$. To check the numerical hypothesis of the connector
lemma, use $g\lfloor R/g\rfloor\ge R-g$ and $g\le k$. These give
\[
 2g\lfloor R/g\rfloor
 \ge2R-2g
 \ge\frac{2(k-1)M+2}{k}
 \ge M+k=h.
\]
For the last inequality, it suffices that $(k-2)M+2\ge k^2$.
This follows from $M\ge2k$ and $k\ge5$, since
$2k(k-2)+2\ge k^2$. Thus the connector lemma supplies an actual list
of at most $2g^2\le2(r-1)^2$ circles for each contour.
Summing over the contours gives
$c_{\mathrm{adv}}\le(r-1)q/M\le(r-1)q/m$.

Before deletion the number of candidates is
$C_0=H(m)_G\lfloor\eta M^d\rfloor$. Define
\[
 \phi=\frac{(m)_G}{m^G},\qquad
 \psi=\frac{\lfloor\eta M^d\rfloor}{\eta M^d},\qquad
 \mu=\frac{M}{m},\qquad
 \chi=\frac{D_{m,W}}{m^W}.
\]
The exact leading contributions before deletion are
\begin{align*}
 C_0(r-2)M^2
   &=H\beta m^{W-1}\phi\psi\mu^{d+2},\\
 C_0q_0
   &=c_rH\beta m^{W-2}\phi\psi\mu^{d+1}.
\end{align*}
The exponents follow from $G+d+2=W-1$ and $G+d+1=W-2$.
Let $t$ be the proportion of candidates that survive both deletions.
The deletion estimates give
\[
 1-\frac{18r^2W^3+7\,118\,280r^{22}W^4}{m}\le t\le1.
\]
In particular $t>0$ under the present size bound. Averaging the relative
row errors over the surviving contours gives a number $u$ such that
\[
 |u|\le\frac{2r+1040r^5W^2}{m}.
\]
Their short charges are all exactly $q_0$. Write
$v=c_{\mathrm{adv}}/q$, so $0\le v\le(r-1)/m$. The two normalized
counts now have the exact expressions
\begin{align*}
 \frac{B}{D_{m,W}}
   &=\frac{H\beta}{m}
     \frac{\phi\psi t\mu^{d+2}(1+u)}{\chi},\\
 \frac{q+c_{\mathrm{adv}}}{D_{m,W}}
   &=\frac{c_rH\beta}{m^2}
     \frac{\phi\psi t\mu^{d+1}(1+v)}{\chi}.
\end{align*}
The relative-error numerators, each to be divided by $m$, are bounded by
\[
\begin{array}{lr}
\text{rounding the selected frame count}&W\,,\\
\text{distinct guard positions}&W^2/2\,,\\
\text{guard-only deletion}&18r^2W^3\,,\\
\text{nearby incompatible events}&7\,118\,280r^{22}W^4\,,\\
\text{row count within a contour}&2r+1040r^5W^2\,,\\
\text{additional connector charge}&r\,,\\
\text{replacing powers of }M\text{ by powers of }m&2W^2\,,\\
\text{the product }D_{m,W}/m^W&W^2\,.
\end{array}
\]
The first bound uses $d\ge3$ and $\eta^{-1}\le W/r^4$.
The next three were proved in \eqref{eq-guard-count},
\eqref{eq-guard-loss}, and \eqref{eq-proximity-loss}.
The row bound follows by dividing the displayed estimate for $N$ by
$(r-2)M^2$. We use the product inequalities
\[
 \prod_i(1+u_i)\le\frac1{1-\sum_i u_i},\qquad
 \prod_i(1-u_i)\ge1-\sum_i u_i
 \quad(0\le u_i,\ \sum_i u_i<1).
\]
For the first inequality, use $1+u_i\le1/(1-u_i)$ and the second
inequality. For the powers of $\mu$, there are at most $r$ equal factors
$1+G/m$. Their deviations sum to at most $rG/m\le1/4$, so
\[
 1\le\mu^{d+1}\le\mu^{d+2}
 \le\frac1{1-rG/m}\le1+\frac{2rG}{m}
 \le1+\frac{2W^2}{m}.
\]
For $\chi$, the deviations $i/m$ sum to at most $W^2/(2m)\le1/8$,
which gives $1\le\chi\le1+W^2/m$ in the same way.

The sum of the listed numerators is less than $9\cdot10^6r^{22}W^4$.
Let $s$ be that sum divided by $m$. Then $s<1/1000$. To combine the
estimates, let $s_0$ be the sum of the deviations allowed for the
numerator factors in either exact expression, and let $s_1$ be the
deviation allowed for $\chi$. Thus $s_0+s_1\le s$.
The product inequalities give a numerator between $1-s_0$ and
$1/(1-s_0)$, while $1\le\chi\le1+s_1$. The ratio is therefore at
least
\[
 \frac{1-s_0}{1+s_1}\ge1-s_0-s_1\ge1-s
\]
and at most $1/(1-s_0)\le1/(1-s)\le1+2s$.
In particular its deviation from one is at most $4s$.
Since $36\cdot10^6<10^{12}$, each relative error is at most $E/m$.
Finally $H\beta=\rho W$, as required.
\end{proof}

\begin{lemma}\label{lem-polynomial-width}
The parameters in \eqref{eq-global-parameters} satisfy
$r^2\le W\le352r^9$ and $r^2\le\rho<r^2+1$.
\end{lemma}
\begin{proof}
Cancellation of factorials gives
\[
 \beta=\frac{r(r-1)(r-2)^2}{2(64r^8+1)}.
\]
For $r\ge6$, the inequalities $r-1\ge5r/6$ and
$r-2\ge2r/3$ imply $\beta\ge1/(351r^4)$.
Also $0<\beta<1$. Thus
\[
 r^2\le W=r^2A<\frac{r^5}{\beta}+r^2
 \le351r^9+r^2\le352r^9.
\]
The defining ceiling for $A$ gives
$r^3\le A\beta<r^3+\beta$. Division by $r$ yields the bounds on
$\rho$.
\end{proof}
